% Continuación para colocar después de:
% \subsection{Aplicación en ventanas existentes}

Después de comprobar el funcionamiento de cada tecnología con los botones de prueba, pasaremos a utilizarlas desde las opciones que ya tenemos en nuestra ventana de bienvenida. De esta manera, cada botón tendrá una función dentro de la práctica y podremos observar la comunicación entre el backend y Flutter desde una pantalla específica.

La distribución de las ventanas quedará de la siguiente manera:
\begin{itemize}
    \item \textbf{Mi cuenta:} permitirá consultar el perfil y actualizar el teléfono mediante GraphQL.
    \item \textbf{Notificaciones:} recibirá los avisos enviados por el servidor mediante SSE.
    \item \textbf{Chat:} permitirá enviar mensajes al servidor y recibir sus respuestas mediante WebSocket.
    \item \textbf{Privacidad:} mostrará un aviso genérico sobre el uso académico de la aplicación.
\end{itemize}

NOTA: la opción \textbf{Cuenta} que se encuentra dentro del registro seguirá utilizando su formulario original. La nueva ventana \textbf{Mi cuenta} consultará los datos simulados del backend; por lo tanto, no mostrará automáticamente los datos escritos en el formulario de registro ni el perfil del usuario que inició sesión.

\subsubsection{Revisión de las dependencias utilizadas}
Antes de continuar, debemos revisar el archivo \texttt{pubspec.yaml}. El código de SSE mostrado en esta guía utiliza el paquete \texttt{flutter\_client\_sse}, aunque en la lista inicial de dependencias se indicó \texttt{server\_event\_source}. Para trabajar con los archivos que se presentan a continuación, ejecutaremos los siguientes comandos desde la raíz del proyecto Flutter:

\begin{lstlisting}[language=bash, caption={Dependencias utilizadas en las nuevas ventanas}]
flutter pub add flutter_client_sse
flutter pub add http
flutter pub add web_socket_channel
flutter pub add graphql_flutter
flutter pub add flutter_svg
flutter pub get
\end{lstlisting}

Si los paquetes ya se encuentran instalados, no es necesario volver a agregarlos. El paquete \texttt{server\_event\_source} no se utiliza en los archivos de esta implementación.

\subsubsection{Cambio de nombres en el menú principal}
Empezaremos editando el archivo \texttt{assets/es.json}. Buscaremos las propiedades \texttt{products}, \texttt{chatbot} y \texttt{privacy}, y reemplazaremos únicamente sus valores:

\begin{lstlisting}[language=json, caption={Nuevos nombres en assets/es.json}]
"products": "Notificaciones",
"chatbot": "Chat",
"privacy": "Privacidad",
\end{lstlisting}

Debemos hacer el mismo cambio en \texttt{assets/en.json}:

\begin{lstlisting}[language=json, caption={Nuevos nombres en assets/en.json}]
"products": "Notifications",
"chatbot": "Chat",
"privacy": "Privacy",
\end{lstlisting}

NOTA: estos fragmentos son propiedades del JSON existente, no archivos completos. Debemos conservar las llaves y las demás propiedades, sin duplicar las claves. También conservaremos los nombres internos \texttt{products}, \texttt{chatbot} y \texttt{privacy}, ya que el menú los utiliza para identificar cada opción. La aplicación continúa cargando español de forma fija; actualizar el archivo en inglés no agrega un selector de idioma.

\subsubsection{Retirar el botón temporal de SSE}
Ahora volveremos al archivo \texttt{lib/screens/welcome\_page.dart}. Como ya hicimos con GraphQL y WebSocket, comentaremos el botón \textbf{Probar SSE} junto con el espacio que lo acompaña. No debemos borrar los botones del menú principal.

\begin{lstlisting}[language=Java, caption={Botón temporal de SSE comentado}]
/* Botón de prueba SSE oculto para integrar la
   tecnología en una pantalla funcional.
const SizedBox(height: 12.0),

ElevatedButton.icon(
  onPressed: () {
    Navigator.push(
      context,
      MaterialPageRoute(
        builder: (context) => const SseTestPage(),
      ),
    );
  },
  icon: const Icon(Icons.notifications_active),
  label: const Text('Probar SSE'),
),
*/
\end{lstlisting}

Los botones de prueba para GraphQL y WebSocket deben permanecer dentro de su bloque de comentarios. Además, comentaremos las tres importaciones de prueba que se encuentran al inicio de \texttt{welcome\_page.dart}, ya que esta ventana dejará de utilizarlas directamente:

\begin{lstlisting}[language=Java, caption={Importaciones de prueba comentadas}]
// import 'package:aplicacion2/screens/graphql_test_page.dart';
// import 'package:aplicacion2/screens/websocket_test_page.dart';
// import 'package:aplicacion2/screens/sse_test_page.dart';
\end{lstlisting}

Así conservamos el código de las pruebas y dejamos de mostrar sus botones en la parte inferior de la bienvenida. El acceso al chat se realizará desde el menú, por lo que su importación se colocará posteriormente en \texttt{custom\_widgets.dart}.

\subsubsection{Creación del diseño compartido para las ventanas}
Para dar el mismo formato a ChatBot, Notificaciones y Privacidad, crearemos primero el archivo \texttt{lib/widgets/academic\_screen.dart}. En este archivo colocaremos un widget llamado \texttt{AcademicScreen}, que nos permitirá reutilizar el fondo guinda, el logotipo superior, el contenedor blanco y el botón de regreso.

El diseño toma como referencia las pantallas escritas en Kotlin, pero se implementará con widgets de Flutter. Utilizaremos el color guinda \texttt{0xFF6F263B}, un fondo beige para el botón de regreso y un radio de 24 para las esquinas del panel blanco. Las funciones de conexión se conservarán en sus pantallas correspondientes.

Antes de copiar el código, verificaremos que existan los archivos \texttt{assets/logo.svg} y \texttt{assets/Pantera.svg}, respetando las mayúsculas de sus nombres. Estos recursos ya forman parte del proyecto y deben continuar declarados en \texttt{pubspec.yaml}. La importación de \texttt{flutter\_svg} permitirá mostrarlos en Flutter.

Colocaremos el siguiente código completo en el nuevo archivo:

\begin{lstlisting}[language=Java, caption={Código completo de academic\_screen.dart}]
import 'package:flutter/material.dart';
import 'package:flutter_svg/flutter_svg.dart';

const academicGuinda = Color(0xFF6F263B);
const academicBeige = Color(0xFFEBE3D5);

/// Marco visual compartido por las pantallas de la práctica.
class AcademicScreen extends StatelessWidget {
  const AcademicScreen({
    super.key,
    required this.title,
    required this.child,
    this.header,
  });

  final String title;
  final Widget child;
  final Widget? header;

  @override
  Widget build(BuildContext context) {
    return Scaffold(
      backgroundColor: academicGuinda,
      body: SafeArea(
        child: Column(
          children: [
            Padding(
              padding: const EdgeInsets.symmetric(vertical: 16),
              child: SvgPicture.asset(
                'assets/logo.svg',
                height: 40,
                semanticsLabel: 'Logo',
                colorFilter:
                    const ColorFilter.mode(Colors.white, BlendMode.srcIn),
              ),
            ),
            Expanded(
              child: Padding(
                padding: const EdgeInsets.fromLTRB(16, 0, 16, 16),
                child: Material(
                  color: Colors.white,
                  borderRadius: BorderRadius.circular(24),
                  clipBehavior: Clip.antiAlias,
                  child: Column(
                    crossAxisAlignment: CrossAxisAlignment.stretch,
                    children: [
                      Padding(
                        padding: const EdgeInsets.all(16),
                        child: Stack(
                          alignment: Alignment.topCenter,
                          children: [
                            Padding(
                              padding:
                                  const EdgeInsets.symmetric(horizontal: 52),
                              child: header ??
                                  Padding(
                                    padding: const EdgeInsets.symmetric(
                                        vertical: 10),
                                    child: Text(title,
                                        textAlign: TextAlign.center,
                                        style: const TextStyle(
                                            fontSize: 24,
                                            fontWeight: FontWeight.bold,
                                            color: academicGuinda)),
                                  ),
                            ),
                            Align(
                              alignment: Alignment.topLeft,
                              child: IconButton(
                                tooltip: 'Regresar',
                                onPressed: () =>
                                    Navigator.of(context).maybePop(),
                                style: IconButton.styleFrom(
                                  backgroundColor: academicBeige,
                                  foregroundColor: academicGuinda,
                                  shape: RoundedRectangleBorder(
                                      borderRadius: BorderRadius.circular(8)),
                                ),
                                icon: const Icon(Icons.arrow_back),
                              ),
                            ),
                          ],
                        ),
                      ),
                      Expanded(child: child),
                    ],
                  ),
                ),
              ),
            ),
          ],
        ),
      ),
    );
  }
}
\end{lstlisting}

El parámetro \texttt{title} recibe el título de la ventana y \texttt{child} recibe su contenido. El parámetro opcional \texttt{header} permite reemplazar el título por una cabecera personalizada, como la que utilizaremos para mostrar la pantera y el estado de conexión del ChatBot.

El widget utiliza \texttt{SafeArea} para respetar las zonas del sistema y \texttt{Expanded} para distribuir el espacio disponible. El botón \textbf{Regresar} ejecuta \texttt{Navigator.of(context).maybePop()}, por lo que vuelve a la pantalla anterior sin agregar otra bienvenida a la navegación.

NOTA: debemos crear este archivo antes de agregar las nuevas versiones de las tres pantallas, ya que todas importarán \texttt{../widgets/academic\_screen.dart}. La pantalla \textbf{Mi cuenta} mantiene su presentación anterior; el archivo Kotlin de esa ventana no se aplica a las tres pantallas de este cambio estético.

\subsubsection{Creación de la ventana Mi cuenta con GraphQL}
Dentro de \texttt{lib/screens/}, crearemos un archivo llamado \texttt{profile\_page.dart}. Esta ventana utilizará el servicio \texttt{GraphQLService} que creamos anteriormente, por lo que no es necesario volver a escribir las consultas ni modificar \texttt{graphql\_service.dart}.

Al abrir la pantalla, el método \texttt{initState()} ejecutará la consulta del perfil. En la ventana veremos el nombre, el correo electrónico y un campo para editar el teléfono. El botón \textbf{Guardar teléfono} enviará el nuevo valor mediante la mutación de GraphQL.

Para que sea más sencilla la implementación, colocaremos el siguiente código completo:

\begin{lstlisting}[language=Java, caption={Código completo de profile\_page.dart}]
import 'package:flutter/material.dart';
import '../services/graphql_service.dart';

/// Consulta y edición del perfil remoto mediante GraphQL.
class ProfilePage extends StatefulWidget {
  const ProfilePage({super.key});

  @override
  State<ProfilePage> createState() => _ProfilePageState();
}

class _ProfilePageState extends State<ProfilePage> {
  final _formKey = GlobalKey<FormState>();
  final _phone = TextEditingController();
  Map<String, dynamic>? _profile;
  bool _busy = true;
  String? _error;

  @override
  void initState() {
    super.initState();
    _load();
  }

  Future<void> _load() async {
    setState(() {
      _busy = true;
      _error = null;
    });
    try {
      final profile = await GraphQLService.obtenerPerfil();
      if (!mounted) return;
      setState(() {
        _profile = profile;
        _phone.text = profile['telefono']?.toString() ?? '';
      });
    } catch (_) {
      if (mounted) {
        setState(() => _error =
            'No se pudo cargar el perfil. Verifica la conexión con el servidor.');
      }
    } finally {
      if (mounted) setState(() => _busy = false);
    }
  }

  Future<void> _save() async {
    if (!_formKey.currentState!.validate()) return;
    setState(() {
      _busy = true;
      _error = null;
    });
    try {
      final profile =
          await GraphQLService.actualizarTelefono(_phone.text.trim());
      if (!mounted) return;
      setState(() {
        _profile = profile;
        _phone.text = profile['telefono']?.toString() ?? '';
      });
      ScaffoldMessenger.of(context).showSnackBar(
        const SnackBar(content: Text('Teléfono actualizado')),
      );
    } catch (_) {
      if (mounted) {
        setState(() =>
            _error = 'No se pudo actualizar el teléfono. Inténtalo de nuevo.');
      }
    } finally {
      if (mounted) setState(() => _busy = false);
    }
  }

  @override
  void dispose() {
    _phone.dispose();
    super.dispose();
  }

  @override
  Widget build(BuildContext context) {
    return Scaffold(
      appBar: AppBar(title: const Text('Mi cuenta'), actions: [
        IconButton(
            onPressed: _busy ? null : _load,
            tooltip: 'Actualizar perfil',
            icon: const Icon(Icons.refresh)),
      ]),
      body: SingleChildScrollView(
        padding: const EdgeInsets.all(24),
        child: Form(
          key: _formKey,
          child:
              Column(crossAxisAlignment: CrossAxisAlignment.stretch, children: [
            if (_busy) const LinearProgressIndicator(),
            if (_error != null) ...[
              Text(_error!),
              TextButton(
                  onPressed: _busy ? null : _load,
                  child: const Text('Volver a consultar')),
            ],
            if (_profile != null) ...[
              ListTile(
                  leading: const Icon(Icons.person),
                  title: const Text('Nombre'),
                  subtitle:
                      Text(_profile!['nombre']?.toString() ?? 'Sin nombre')),
              ListTile(
                  leading: const Icon(Icons.email),
                  title: const Text('Correo electrónico'),
                  subtitle:
                      Text(_profile!['email']?.toString() ?? 'Sin correo')),
              const SizedBox(height: 16),
              TextFormField(
                  controller: _phone,
                  enabled: !_busy,
                  keyboardType: TextInputType.phone,
                  decoration: const InputDecoration(
                      labelText: 'Teléfono', border: OutlineInputBorder()),
                  validator: (value) {
                    final phone = value?.trim() ?? '';
                    if (!RegExp(r'^\+?[0-9 ()-]{7,20}$').hasMatch(phone) ||
                        phone.replaceAll(RegExp(r'\D'), '').length < 7) {
                      return 'Ingresa un teléfono válido.';
                    }
                    return null;
                  }),
              const SizedBox(height: 16),
              ElevatedButton(
                  onPressed: _busy ? null : _save,
                  child: const Text('Guardar teléfono')),
            ],
          ]),
        ),
      ),
    );
  }
}
\end{lstlisting}

La variable \texttt{\_busy} permite mostrar el indicador de carga y deshabilitar los controles mientras se procesa una solicitud. Si ocurre un error, aparecerá un mensaje y podremos volver a consultar el perfil. Antes de enviar el teléfono, se validará que tenga un formato permitido y por lo menos siete dígitos.

NOTA: el servicio existente utiliza la operación \texttt{actualizarTelefono}. Este es el nombre utilizado por la consulta de Flutter para la mutación del backend mostrado anteriormente. Debemos conservarlo al copiar el código del servicio.

\subsubsection{Creación de la ventana Notificaciones con SSE}
El siguiente paso será crear el archivo \texttt{lib/screens/notifications\_page.dart}. Podemos tomar como base la pantalla temporal \texttt{sse\_test\_page.dart}, pero debemos conservar ambos archivos y cambiar el nombre de la nueva clase a \texttt{NotificationsPage}.

En esta versión, la suscripción se iniciará al abrir la ventana. También tendremos los botones \textbf{Pausar notificaciones} y \textbf{Recibir notificaciones}, que permitirán detener y reactivar la recepción sin salir de la pantalla. Conservaremos el botón \textbf{Generar notificación} para demostrar el funcionamiento durante la práctica.

La ventana utilizará el fondo guinda y el panel blanco definidos en \texttt{AcademicScreen}. Dentro del panel veremos el título centrado, el estado de recepción y los controles de la práctica. Las notificaciones aparecerán en tarjetas grises y, cuando la lista esté vacía, se mostrará el mensaje \textit{No hay notificaciones nuevas}. Los controles se colocarán dentro de un \texttt{Wrap} para distribuirse en varias filas cuando sea necesario.

Colocaremos el siguiente código completo en el nuevo archivo:

\begin{lstlisting}[language=Java, caption={Código completo de notifications\_page.dart}]
import 'dart:async';
import 'dart:convert';

import 'package:flutter/material.dart';
import 'package:flutter_client_sse/flutter_client_sse.dart';
import 'package:flutter_client_sse/constants/sse_request_type_enum.dart';
import 'package:http/http.dart' as http;
import '../widgets/academic_screen.dart';

class NotificationsPage extends StatefulWidget {
  const NotificationsPage({super.key});

  @override
  State<NotificationsPage> createState() => _NotificationsPageState();
}

class _NotificationsPageState extends State<NotificationsPage> {
  StreamSubscription<dynamic>? _subscription;

  final List<String> _notificaciones = [];

  bool _conectado = false;
  bool _generando = false;
  String _estado = 'Desconectado';

  @override
  void initState() {
    super.initState();
    _conectar();
  }

  void _conectar() {
    if (_conectado) return;

    _subscription?.cancel();
    SSEClient.unsubscribeFromSSE();

    setState(() {
      _estado = 'Conectando al servidor...';
    });

    try {
      _subscription = SSEClient.subscribeToSSE(
        method: SSERequestType.GET,
        url: 'http://localhost:8000/notificaciones/eventos',
        header: {
          'Accept': 'text/event-stream',
          'Cache-Control': 'no-cache',
        },
      ).listen(
        (evento) {
          if (!mounted) return;
          if (evento.data == null || evento.data!.trim().isEmpty) return;

          setState(() {
            _conectado = true;
            _estado = 'Conectado: escuchando notificaciones';

            try {
              final datos = jsonDecode(evento.data ?? '{}');

              _notificaciones.insert(
                0,
                'ID: ${datos['id']}\n'
                'Mensaje: ${datos['mensaje']}\n'
                'Hora: ${datos['hora']}',
              );
            } catch (_) {
              _notificaciones.insert(
                0,
                'Evento recibido: ${evento.data ?? ''}',
              );
            }
          });
        },
        onError: (Object error) {
          if (!mounted) return;

          setState(() {
            _conectado = false;
            _estado = 'Error SSE: $error';
          });
        },
        onDone: () {
          if (!mounted) return;

          setState(() {
            _conectado = false;
            _estado = 'Conexión SSE cerrada';
          });
        },
      );

      setState(() {
        _conectado = true;
        _estado = 'Suscripción iniciada; esperando eventos...';
      });
    } catch (error) {
      setState(() {
        _conectado = false;
        _estado = 'No se pudo conectar: $error';
      });
    }
  }

  Future<void> _generarNotificacion() async {
    if (_generando) return;

    setState(() {
      _generando = true;
    });

    try {
      final respuesta = await http
          .post(
            Uri.parse('http://localhost:8000/notificaciones/generar'),
            headers: {
              'Content-Type': 'application/json',
            },
            body: jsonEncode({
              'mensaje': 'Notificación de prueba desde Flutter',
            }),
          )
          .timeout(const Duration(seconds: 15));

      if (!mounted) return;

      if (respuesta.statusCode >= 200 && respuesta.statusCode < 300) {
        setState(() {
          _estado = 'Notificación generada correctamente';
        });
      } else {
        setState(() {
          _estado = 'Error HTTP ${respuesta.statusCode}: ${respuesta.body}';
        });
      }
    } catch (error) {
      if (!mounted) return;

      setState(() {
        _estado = 'Error al generar notificación: $error';
      });
    } finally {
      if (mounted) {
        setState(() {
          _generando = false;
        });
      }
    }
  }

  void _desconectar() {
    _subscription?.cancel();
    _subscription = null;

    SSEClient.unsubscribeFromSSE();

    setState(() {
      _conectado = false;
      _estado = 'Desconectado';
    });
  }

  @override
  void dispose() {
    _subscription?.cancel();
    SSEClient.unsubscribeFromSSE();
    super.dispose();
  }

  @override
  Widget build(BuildContext context) {
    return AcademicScreen(
      title: 'Notificaciones',
      child: Padding(
        padding: const EdgeInsets.fromLTRB(16, 0, 16, 16),
        child: Column(
          crossAxisAlignment: CrossAxisAlignment.stretch,
          children: [
            Text(
              _estado,
              textAlign: TextAlign.center,
              style: TextStyle(
                fontWeight: FontWeight.bold,
                color: _conectado ? Colors.green : Colors.red,
              ),
            ),
            const SizedBox(height: 16),
            Wrap(
              alignment: WrapAlignment.center,
              spacing: 8,
              runSpacing: 8,
              children: [
                OutlinedButton(
                  style:
                      OutlinedButton.styleFrom(foregroundColor: academicGuinda),
                  onPressed: _conectado ? null : _conectar,
                  child: const Text('Recibir notificaciones'),
                ),
                OutlinedButton(
                  style:
                      OutlinedButton.styleFrom(foregroundColor: academicGuinda),
                  onPressed: _conectado ? _desconectar : null,
                  child: const Text('Pausar notificaciones'),
                ),
                ElevatedButton.icon(
                  style: ElevatedButton.styleFrom(
                      backgroundColor: academicGuinda,
                      foregroundColor: Colors.white,
                      shape: RoundedRectangleBorder(
                          borderRadius: BorderRadius.circular(8))),
                  onPressed: _generando ? null : _generarNotificacion,
                  icon: const Icon(Icons.notifications_active),
                  label: Text(
                    _generando ? 'Generando...' : 'Generar notificación',
                  ),
                ),
              ],
            ),
            const Divider(height: 32),
            const Text(
              'Notificaciones recibidas',
              style: TextStyle(
                fontSize: 18,
                fontWeight: FontWeight.bold,
                color: academicGuinda,
              ),
            ),
            const SizedBox(height: 8),
            Expanded(
              child: _notificaciones.isEmpty
                  ? const Center(
                      child: Text(
                        'No hay notificaciones nuevas',
                        textAlign: TextAlign.center,
                        style: TextStyle(color: Colors.grey),
                      ),
                    )
                  : ListView.builder(
                      itemCount: _notificaciones.length,
                      itemBuilder: (context, index) {
                        return Card(
                          color: const Color(0xFFF5F5F5),
                          surfaceTintColor: Colors.transparent,
                          elevation: 0,
                          margin: const EdgeInsets.symmetric(vertical: 8),
                          child: ListTile(
                            leading: const Icon(
                              Icons.notifications,
                              color: academicGuinda,
                            ),
                            title: Text(
                              _notificaciones[index],
                              style: const TextStyle(color: Colors.black87),
                            ),
                          ),
                        );
                      },
                    ),
            ),
          ],
        ),
      ),
    );
  }
}
\end{lstlisting}

Con estos cambios, el método \texttt{\_conectar()} se ejecutará desde \texttt{initState()}. Antes de iniciar una nueva suscripción se cancelará la anterior, y al salir de la ventana también se liberará la conexión. Los eventos sin datos no se agregarán a la lista.

El botón \textbf{Generar notificación} realizará una petición HTTP al endpoint \texttt{/notificaciones/generar}. El backend enviará el aviso a las conexiones SSE activas y la pantalla lo agregará a la lista cuando reciba el evento. La respuesta de la petición HTTP no inserta directamente una notificación en la ventana.

NOTA: el texto \textit{Suscripción iniciada; esperando eventos...} indica que se inició la suscripción desde Flutter; la recepción de una notificación permite comprobar el recorrido completo del evento. La petición para generar un aviso tiene un tiempo de espera de 15 segundos y acepta respuestas HTTP del rango 200 al 299.

\subsubsection{Adaptación de la ventana WebSocket al diseño de ChatBot}
Para el chat reutilizaremos el archivo \texttt{lib/screens/websocket\_test\_page.dart}. Conservaremos el nombre de la clase \texttt{WebSocketTestPage}, porque el menú ya utiliza esa clase para abrir la ventana. Reemplazaremos su contenido por el siguiente código completo, que incorpora el diseño nuevo y conserva la conexión WebSocket existente:

\begin{lstlisting}[language=Java, caption={Código completo de websocket\_test\_page.dart}]
import 'dart:async';

import 'package:flutter/material.dart';
import 'package:web_socket_channel/web_socket_channel.dart';
import 'package:flutter_svg/flutter_svg.dart';
import '../widgets/academic_screen.dart';

class WebSocketTestPage extends StatefulWidget {
  const WebSocketTestPage({super.key});

  @override
  State<WebSocketTestPage> createState() => _WebSocketTestPageState();
}

class _WebSocketTestPageState extends State<WebSocketTestPage> {
  WebSocketChannel? _channel;
  StreamSubscription<dynamic>? _subscription;

  final TextEditingController _messageController = TextEditingController();

  final List<String> _messages = [];

  bool _connected = false;
  String _status = 'Conectando al servidor...';

  @override
  void initState() {
    super.initState();
    _connect();
  }

  void _connect() {
    try {
      final channel = WebSocketChannel.connect(
        Uri.parse('ws://localhost:8000/ws/chat'),
      );

      _channel = channel;

      channel.ready.then<void>((_) {
        if (!mounted) return;

        setState(() {
          _connected = true;
          _status = 'Conectado al servidor WebSocket.';
        });
      }).catchError((Object error) {
        if (!mounted) return;

        setState(() {
          _connected = false;
          _status = 'Error de conexión: $error';
        });
      });

      _subscription = channel.stream.listen(
        (mensaje) {
          if (!mounted) return;

          setState(() {
            _messages.add('Servidor: $mensaje');
          });
        },
        onError: (Object error) {
          if (!mounted) return;

          setState(() {
            _connected = false;
            _status = 'Error de WebSocket: $error';
          });
        },
        onDone: () {
          if (!mounted) return;

          setState(() {
            _connected = false;
            _status = 'Conexión cerrada.';
          });
        },
      );
    } catch (error) {
      setState(() {
        _connected = false;
        _status = 'No se pudo conectar: $error';
      });
    }
  }

  void _sendMessage() {
    final mensaje = _messageController.text.trim();

    if (mensaje.isEmpty || !_connected || _channel == null) {
      return;
    }

    _channel!.sink.add(mensaje);

    setState(() {
      _messages.add('Tú: $mensaje');
    });

    _messageController.clear();
  }

  @override
  void dispose() {
    _subscription?.cancel();
    _channel?.sink.close();
    _messageController.dispose();

    super.dispose();
  }

  @override
  Widget build(BuildContext context) {
    return AcademicScreen(
      title: 'ChatBot',
      header: Column(
        children: [
          SvgPicture.asset('assets/Pantera.svg',
              height: 50,
              width: 50,
              semanticsLabel: 'Icono ChatBot',
              colorFilter:
                  const ColorFilter.mode(academicGuinda, BlendMode.srcIn)),
          const SizedBox(height: 4),
          const Text('ChatBot',
              style: TextStyle(
                  fontSize: 20,
                  fontWeight: FontWeight.bold,
                  color: academicGuinda)),
          const Text('Tu asistente virtual',
              style: TextStyle(fontSize: 12, color: Colors.grey)),
          const SizedBox(height: 4),
          Tooltip(
              message: _status,
              child: Text(_status,
                  maxLines: 2,
                  overflow: TextOverflow.ellipsis,
                  textAlign: TextAlign.center,
                  style: TextStyle(
                      fontSize: 12,
                      color: _connected ? Colors.green : Colors.red))),
        ],
      ),
      child: Padding(
        padding: const EdgeInsets.fromLTRB(16, 0, 16, 16),
        child: Column(
          children: [
            Expanded(
              child: _messages.isEmpty
                  ? const Center(
                      child: Text(
                          'Envía un mensaje para conversar con el servidor.',
                          textAlign: TextAlign.center,
                          style: TextStyle(color: Colors.grey, fontSize: 14)))
                  : ListView.builder(
                      itemCount: _messages.length,
                      itemBuilder: (context, index) {
                        final isMine = _messages[index].startsWith('Tú:');
                        return Padding(
                          padding: const EdgeInsets.symmetric(vertical: 6),
                          child: Text(_messages[index],
                              style: TextStyle(
                                  color: isMine
                                      ? academicGuinda
                                      : Colors.black87)),
                        );
                      },
                    ),
            ),
            Row(
              children: [
                Expanded(
                  child: TextField(
                    controller: _messageController,
                    enabled: _connected,
                    decoration: const InputDecoration(
                      hintText: 'Escribe tu mensaje...',
                      contentPadding:
                          EdgeInsets.symmetric(horizontal: 16, vertical: 14),
                      border: OutlineInputBorder(
                          borderRadius: BorderRadius.all(Radius.circular(24)),
                          borderSide: BorderSide(color: Colors.grey)),
                      focusedBorder: OutlineInputBorder(
                          borderRadius: BorderRadius.all(Radius.circular(24)),
                          borderSide: BorderSide(color: academicGuinda)),
                    ),
                    onSubmitted: (_) => _sendMessage(),
                  ),
                ),
                const SizedBox(width: 8),
                IconButton(
                  onPressed: _connected ? _sendMessage : null,
                  icon: const Icon(Icons.send),
                  tooltip: 'Enviar mensaje',
                  style: IconButton.styleFrom(
                    minimumSize: const Size(50, 50),
                    backgroundColor: academicGuinda,
                    foregroundColor: Colors.white,
                    disabledBackgroundColor: academicBeige,
                    shape: const CircleBorder(),
                  ),
                ),
              ],
            ),
          ],
        ),
      ),
    );
  }
}
\end{lstlisting}

En esta versión, \texttt{AcademicScreen} mostrará el fondo guinda, el logo y el panel blanco. En el parámetro \texttt{header} colocaremos la imagen de la pantera, el título \textbf{ChatBot}, el texto \textit{Tu asistente virtual} y el estado real de la conexión. Los mensajes propios se mostrarán en guinda y los mensajes recibidos en gris oscuro.

En la parte inferior tendremos el campo redondeado con el texto \textit{Escribe tu mensaje...} y un botón circular guinda para enviar. El envío seguirá disponible mediante ese botón o con Enter. Si la conexión no está activa, el campo y el botón permanecerán deshabilitados, como en la versión anterior.

NOTA: el título \textbf{ChatBot} corresponde a la presentación de la ventana. No estamos agregando un modelo de inteligencia artificial ni cambiando las respuestas del backend. La dirección continúa siendo \texttt{ws://localhost:8000/ws/chat}; no debemos copiar la dirección ni el puerto del ejemplo Kotlin.

Con el archivo \texttt{main.py} presentado en esta guía, el servidor enviará primero el mensaje \texttt{bienvenido} y responderá \texttt{hola} cada vez que reciba un texto. Esta implementación demuestra comunicación entre Flutter y el servidor; todavía no distribuye los mensajes de un usuario a los demás clientes. No debemos esperar que dos ventanas del navegador conversen entre sí con ese backend.

\subsubsection{Creación de la ventana Privacidad}
Ahora crearemos el archivo \texttt{lib/screens/privacy\_page.dart}. Esta ventana mostrará un aviso genérico para el uso académico de la aplicación, por lo que no necesita conectarse al backend. Para mantener el mismo formato que ChatBot y Notificaciones, utilizaremos también \texttt{AcademicScreen}, con el título \textbf{Aviso de privacidad}, el panel blanco redondeado y el botón beige de regreso. El contenido permanecerá dentro de un \texttt{SingleChildScrollView} para poder leerlo completo.

\begin{lstlisting}[language=Java, caption={Código completo de privacy\_page.dart}]
import 'package:flutter/material.dart';
import '../widgets/academic_screen.dart';

class PrivacyPage extends StatelessWidget {
  const PrivacyPage({super.key});

  @override
  Widget build(BuildContext context) {
    return const AcademicScreen(
      title: 'Aviso de privacidad',
      child: SingleChildScrollView(
        padding: EdgeInsets.fromLTRB(24, 0, 24, 24),
        child: Text(
          'Proyecto académico\n\n'
          'Esta aplicación se utiliza con fines educativos para demostrar '
          'autenticación, formularios y comunicación con servidores.\n\n'
          'Datos utilizados\n\n'
          'Los formularios permiten introducir datos personales, de cuenta y '
          'domicilio. El inicio de sesión envía las credenciales al servidor '
          'configurado y guarda los tokens de sesión en el dispositivo. '
          'Los mensajes del chat se envían al servidor configurado.\n\n'
          'Uso de la información\n\n'
          'Para las prácticas se recomienda utilizar datos ficticios y '
          'contraseñas exclusivas para este proyecto. La conservación de '
          'información enviada depende del servidor utilizado.\n\n'
          'Control de la sesión\n\n'
          'Al cerrar sesión se eliminan los tokens guardados por la aplicación '
          'y se limpian los datos temporales del formulario. Esto no elimina '
          'información que ya se haya enviado al servidor.\n\n'
          'Este aviso describe el uso académico de la aplicación y deberá '
          'actualizarse si cambian sus funciones o el tratamiento de datos.',
          style: TextStyle(fontSize: 16, height: 1.5, color: Colors.black87),
        ),
      ),
    );
  }
}
\end{lstlisting}

El texto describe los formularios, el inicio de sesión y los mensajes enviados al servidor. También explica que cerrar sesión elimina los tokens locales, pero no borra la información que ya se haya enviado al backend.

\subsubsection{Conexión de las nuevas ventanas con el menú}
Ya que tenemos listas las pantallas, debemos modificar \texttt{lib/widgets/custom\_widgets.dart}. Este archivo contiene la función \texttt{customIconButton}, que determina lo que ocurre al presionar cada opción de la bienvenida.

Retiraremos de este archivo la importación de \texttt{account\_page.dart}, ya que el acceso \textbf{Mi cuenta} abrirá ahora \texttt{ProfilePage}. Esto no elimina el formulario de cuenta dentro del registro: \texttt{register\_page.dart} conserva su propia importación y seguirá utilizándolo.

Las nuevas importaciones y los cambios de navegación quedarán incluidos en el siguiente código completo de \texttt{custom\_widgets.dart}:

\begin{lstlisting}[language=Java, caption={Código completo de custom\_widgets.dart}]
import 'package:flutter/material.dart';
import 'package:flutter_svg/flutter_svg.dart';
import 'package:aplicacion2/services/token_manager.dart';
import 'package:aplicacion2/models/userdata.dart';
import 'package:aplicacion2/screens/register_page.dart';
import 'package:aplicacion2/screens/profile_page.dart';
import 'package:aplicacion2/screens/notifications_page.dart';
import 'package:aplicacion2/screens/websocket_test_page.dart';
import 'package:aplicacion2/screens/privacy_page.dart';

/// Logo fijo y posicionado en la parte superior,
/// centrado horizontalmente.
Widget logo() {
  return Positioned(
    top: 0,
    left: 0,
    right: 0,
    child: SvgPicture.asset(
      'assets/logo.svg',
      height: 130.0,
      width: 130.0,
      colorFilter: const ColorFilter.mode(
        Colors.white,
        BlendMode.srcIn,
      ),
    ),
  );
}

/// Botón personalizado con ícono SVG y texto.
///
/// [nombreImagen] es la ruta del ícono SVG.
/// [acronimo] es el texto mostrado debajo del ícono.
/// [context] es el contexto de la aplicación.
/// [jsonStrings] contiene las traducciones utilizadas
/// para la navegación.
Widget customIconButton({
  required String nombreImagen,
  required String acronimo,
  required BuildContext context,
  required Map<String, dynamic> jsonStrings,
}) {
  return InkWell(
    child: SizedBox(
      width: 130,
      height: 130,
      child: Card(
        margin: const EdgeInsets.all(10.0),
        elevation: 9.0,
        child: Column(
          mainAxisAlignment: MainAxisAlignment.center,
          children: [
            SvgPicture.asset(
              nombreImagen,
              width: 45,
              height: 45,
              colorFilter: const ColorFilter.mode(
                Color.fromARGB(255, 111, 38, 59),
                BlendMode.srcIn,
              ),
            ),
            Text(
              acronimo,
              style: const TextStyle(
                color: Color.fromARGB(255, 111, 38, 59),
                fontSize: 15.0,
              ),
            ),
          ],
        ),
      ),
    ),

    // Acción al presionar el botón.
    onTap: () async {
      debugPrint('Click en $acronimo');

      // Navegación hacia Registro.
      if (acronimo == 'Registro' || acronimo == 'Register') {
        Navigator.push(
          context,
          MaterialPageRoute(
            builder: (context) => RegisterPage(
              jsonStrings: jsonStrings,
            ),
          ),
        );
      }

      // Navegación hacia Mi Cuenta.
      else if (acronimo == jsonStrings['my_account']) {
        Navigator.push(
          context,
          MaterialPageRoute(
            builder: (context) => const ProfilePage(),
          ),
        );
      } else if (acronimo == jsonStrings['products']) {
        Navigator.push(
            context,
            MaterialPageRoute(
              builder: (_) => const NotificationsPage(),
            ));
      } else if (acronimo == jsonStrings['chatbot']) {
        Navigator.push(
            context,
            MaterialPageRoute(
              builder: (_) => const WebSocketTestPage(),
            ));
      } else if (acronimo == jsonStrings['privacy']) {
        Navigator.push(
            context,
            MaterialPageRoute(
              builder: (_) => const PrivacyPage(),
            ));
      }
      // 3. LÓGICA DE CERRAR SESIÓN
      else if (acronimo == jsonStrings['logout'] ||
          acronimo == 'Cerrar sesión') {
        // Borrar los tokens guardados en SharedPreferences.
        await TokenManager.clearTokens();

        // Limpiar los datos de la sesión actual.
        UserData.clear();

        // Regresar al login y eliminar las pantallas anteriores.
        if (!context.mounted) return;

        Navigator.pushNamedAndRemoveUntil(
          context,
          '/login',
          (route) => false,
        );
      }
    },
  );
}

/// Campo de texto reutilizable con validación de error
/// y soporte multilenguaje.
Widget buildField(
  TextEditingController controller,
  String label,
  bool hasError,
  Map<String, dynamic> jsonStrings,
) {
  return Column(
    crossAxisAlignment: CrossAxisAlignment.start,
    children: [
      TextFormField(
        controller: controller,
        style: const TextStyle(fontSize: 14),
        decoration: InputDecoration(
          isDense: true,
          contentPadding: const EdgeInsets.symmetric(
            vertical: 12.0,
            horizontal: 12.0,
          ),
          border: const OutlineInputBorder(),
          labelText: label,
          labelStyle: const TextStyle(fontSize: 14),
        ),
      ),
      if (hasError)
        Align(
          alignment: Alignment.centerLeft,
          child: Text(
            jsonStrings['mandatory_field'],
            style: const TextStyle(
              color: Color.fromARGB(255, 80, 39, 51),
              fontSize: 12,
            ),
          ),
        ),
      const SizedBox(height: 14.0),
    ],
  );
}

/// Widget auxiliar para mostrar un elemento de lista
/// con una viñeta (bullet).
Widget bulletItem(String text) {
  return Padding(
    padding: const EdgeInsets.only(bottom: 0.0),
    child: Row(
      crossAxisAlignment: CrossAxisAlignment.start,
      children: [
        const Text(
          '• ',
          style: TextStyle(
            fontSize: 18.0,
            color: Color.fromARGB(255, 111, 38, 59),
          ),
        ),
        Expanded(
          child: Text(
            text,
            style: const TextStyle(
              color: Color.fromARGB(255, 111, 38, 59),
              fontSize: 14.0,
              fontWeight: FontWeight.bold,
            ),
          ),
        ),
      ],
    ),
  );
}
\end{lstlisting}

Aquí podemos observar que cada opción se compara con su valor en \texttt{jsonStrings}. \textbf{Mi cuenta} abrirá \texttt{ProfilePage}; \textbf{Notificaciones}, \texttt{NotificationsPage}; \textbf{Chat}, \texttt{WebSocketTestPage}; y \textbf{Privacidad}, \texttt{PrivacyPage}. Las acciones de registro y cierre de sesión se conservan. También se utiliza \texttt{debugPrint} para mostrar en consola la opción seleccionada.

\subsubsection{Corrección de la clave de dirección en Registro}
Antes de ejecutar la práctica, revisaremos también el archivo \texttt{lib/screens/register\_page.dart}. Los archivos de idioma utilizan la clave \texttt{adress}, pero la pantalla de registro buscaba \texttt{address} en dos lugares. Esto devuelve un valor nulo y provoca el error \textit{Null is not a subtype of String} al abrir la ventana.

Buscaremos las dos apariciones del siguiente código:

\begin{lstlisting}[language=Java, caption={Clave que debemos corregir en Registro}]
jsonStrings['address'],
\end{lstlisting}

Y las reemplazaremos por:

\begin{lstlisting}[language=Java, caption={Clave utilizada por los archivos de idioma}]
jsonStrings['adress'],
\end{lstlisting}

Una aparición corresponde al botón de la sección Dirección y la otra a su etiqueta. Debemos conservar la misma clave en ambos lugares. Después realizaremos un reinicio de la aplicación para comprobar que el registro abre correctamente.

\subsubsection{Actualización de la prueba de navegación}
Si nuestro proyecto incluye el archivo \texttt{test/navigation\_flow\_test.dart}, debemos actualizar la parte de la prueba que abre \textbf{Mi cuenta} desde la bienvenida. Antes esperaba encontrar el formulario de registro; ahora debe comprobar que se muestra el perfil GraphQL.

Buscaremos el paso 8 de esa prueba y reemplazaremos su comentario y sus comprobaciones por el siguiente fragmento. Las instrucciones que regresan a la bienvenida se conservarán después de este bloque:

\begin{lstlisting}[language=Java, caption={Comprobación del nuevo destino de Mi cuenta}]
// 8. Navegar al perfil GraphQL desde WelcomePage.
await _tapOnText(tester, 'Mi cuenta');
expect(find.text('Mi cuenta'), findsOneWidget);
expect(find.byTooltip('Actualizar perfil'), findsOneWidget);
\end{lstlisting}

NOTA: esta comprobación revisa el destino de la navegación. No demuestra por sí sola que el backend haya respondido correctamente ni que el teléfono se haya guardado; eso se verificará con las pruebas manuales siguientes.

\subsubsection{Revisión de las direcciones de conexión}
Antes de ejecutar el proyecto, revisaremos que las direcciones del cliente coincidan con las del backend. Para la práctica local se utilizan las siguientes:

\begin{itemize}
    \item GraphQL: \url{http://localhost:8000/graphql}.
    \item Recepción SSE: \url{http://localhost:8000/notificaciones/eventos}.
    \item Generación de avisos: \url{http://localhost:8000/notificaciones/generar}.
    \item WebSocket: \texttt{ws://localhost:8000/ws/chat}.
\end{itemize}

Las direcciones se encuentran en \texttt{graphql\_service.dart}, \texttt{notifications\_page.dart} y \texttt{websocket\_test\_page.dart}. Si trabajamos con Edge en el mismo equipo donde corre Python, podremos utilizar \texttt{localhost}. Si utilizamos un teléfono o emulador, debemos cambiarlo por una dirección que permita alcanzar el equipo donde se ejecuta el backend.

NOTA: mantendremos Flutter Web en el puerto 3000, que ya está incluido en la configuración CORS de \texttt{main.py}. Si usamos otro origen para la aplicación, debemos revisar la lista \texttt{allow\_origins} del backend. El servidor de autenticación configurado en \texttt{AuthService} es independiente de este backend de práctica; Python no sustituye al servicio que permite iniciar sesión.

\subsubsection{Ejecución y comprobación de las ventanas}
Primero abriremos una terminal en la carpeta del backend, activaremos el entorno virtual de acuerdo con el sistema utilizado y ejecutaremos:

\begin{lstlisting}[language=bash, caption={Inicio del backend para la práctica integrada}]
python -m uvicorn main:app --reload --host 0.0.0.0 --port 8000
\end{lstlisting}

Esta terminal debe permanecer abierta. En otra terminal, ubicada en la raíz del proyecto Flutter, ejecutaremos:

\begin{lstlisting}[language=bash, caption={Inicio de Flutter Web}]
flutter pub get
flutter run -d edge --web-port 3000
\end{lstlisting}

Después de iniciar sesión veremos el menú con los nombres nuevos y sin los botones temporales en la parte inferior. ChatBot, Notificaciones y Privacidad deberán mostrar el fondo guinda, el logotipo superior y el panel blanco redondeado. En estas tres ventanas el botón beige de regreso debe llevarnos al menú. Para comprobar el funcionamiento realizaremos lo siguiente:

\begin{enumerate}
    \item \textbf{Mi cuenta.} Abriremos esta opción y esperaremos la consulta automática. Deben aparecer el nombre, el correo y el teléfono del perfil simulado en Python. Cambiaremos el teléfono por \texttt{5512345678} y presionaremos \textbf{Guardar teléfono}. Después utilizaremos el icono de actualización para volver a consultar y comprobar que el backend conserva el nuevo valor.

    NOTA: el perfil del backend se guarda en memoria. Si reiniciamos Python, los datos volverán a los valores definidos en \texttt{perfil\_db}.

    \item \textbf{Notificaciones.} Abriremos la ventana y esperaremos el inicio de la suscripción. Presionaremos \textbf{Generar notificación}; al recibir el evento veremos su identificador, mensaje y hora. No es necesario recargar la pantalla.

    \item \textbf{Pausa de notificaciones.} Presionaremos \textbf{Pausar notificaciones} y después \textbf{Generar notificación}. El nuevo aviso no deberá agregarse mientras la recepción esté detenida. Reactivaremos la suscripción con \textbf{Recibir notificaciones}, esperaremos su inicio y generaremos otro aviso. Los avisos anteriores permanecen en la lista mientras no salgamos de la pantalla.

    NOTA: el backend mostrado no guarda un historial para recuperarlo al reconectar. Los avisos generados mientras un cliente está desconectado no se entregan posteriormente a ese cliente. También podemos generar un aviso desde \url{http://localhost:8000/docs} para observar su llegada a Flutter sin presionar el botón de la aplicación.

    \item \textbf{ChatBot.} Abriremos la opción Chat y comprobaremos que aparezcan la pantera, el título y el campo redondeado de mensajes. Después esperaremos el estado de conexión. Debemos recibir \texttt{bienvenido}. Escribiremos un mensaje y lo enviaremos con el botón correspondiente o con Enter; el servidor responderá \texttt{hola}, de acuerdo con el backend de esta guía.

    \item \textbf{Privacidad.} Abriremos el aviso y comprobaremos que el título se muestra dentro del panel blanco, que el texto puede desplazarse y que el botón beige permite regresar al menú.

    \item \textbf{Registro.} Entraremos al registro y abriremos su sección \textbf{Cuenta}. Debe seguir apareciendo el formulario original y no la nueva pantalla del perfil GraphQL.
\end{enumerate}

Para documentar la práctica podemos agregar las siguientes capturas. Debemos sustituir las rutas por las imágenes que hayamos guardado; los bloques permanecen comentados hasta contar con los archivos:

% \begin{figure}[H]
%     \centering
%     \includegraphics[width=0.7\linewidth]{images/perfil_graphql}
%     \caption{Consulta del perfil y actualización del teléfono con GraphQL.}
%     \label{fig:perfil-graphql-integrado}
% \end{figure}

% \begin{figure}[H]
%     \centering
%     \includegraphics[width=0.7\linewidth]{images/notificaciones_sse}
%     \caption{Recepción de notificaciones mediante SSE.}
%     \label{fig:notificaciones-sse-integradas}
% \end{figure}

% \begin{figure}[H]
%     \centering
%     \includegraphics[width=0.7\linewidth]{images/chat_websocket}
%     \caption{Envío de un mensaje y respuesta del servidor mediante WebSocket.}
%     \label{fig:chat-websocket-integrado}
% \end{figure}

\subsubsection{Comprobaciones del código}
Además de las comprobaciones manuales, podemos revisar los archivos nuevos y la navegación con los siguientes comandos desde la raíz de Flutter:

\begin{lstlisting}[language=bash, caption={Análisis de las pantallas integradas}]
flutter analyze --no-pub lib/widgets/academic_screen.dart lib/screens/websocket_test_page.dart lib/screens/notifications_page.dart lib/screens/privacy_page.dart
flutter analyze --no-pub lib/screens/profile_page.dart lib/widgets/custom_widgets.dart
\end{lstlisting}

\begin{lstlisting}[language=bash, caption={Ejecución de las pruebas del proyecto}]
flutter test --no-pub
\end{lstlisting}

El análisis permite detectar problemas del código, pero no comprueba la conexión real con Python. Las pruebas existentes tampoco sustituyen las comprobaciones manuales de los eventos SSE, los mensajes WebSocket y la actualización del teléfono. Si alguna prueba falla, debemos revisar su mensaje y el paso donde ocurre antes de dar por terminada la práctica.

Finalmente, podemos consultar el archivo de apoyo \texttt{docs/comunicacion.md} del proyecto. Este archivo reúne las direcciones utilizadas, el contrato esperado del backend y un guion de demostración. Con los pasos anteriores, las tres tecnologías quedan accesibles desde las opciones del menú principal.
